\begin{ex}
Un intero $x$ si scrive in base $b>1$ se vale l'uguaglianza
\[
x=\sum_{i=0}^{d}\alpha_i b^i
\]
con $\alpha_i\in\{0,1,\dots,b-1\}$ e $\alpha_d\neq 0$; in tal caso, la parola
$[x]_b=\alpha_d\alpha_{d-1}\cdots\alpha_0$ di alfabeto $A=\{0,1,\dots,b-1\}$ si dice la
rappresentazione di $x$ in base $b$.

Dato un insieme $X\subset\mathbb{N}$, denotiamo con
\[
[X]_b=\{[x]_b\mid x\in X\}.
\]
Un insieme $X$ si dice \emph{$b$-automatico} se $[X]_b$ è DFA-decidibile (ovvero regolare).

Mostrare che:
\begin{enumerate}
    \item $X=\mathbb{N}$ è $b$-automatico per ogni intero $b>1$.
    \item $X=\{2^k\mid k\in\mathbb{N}\}$ è $2$-automatico.
    \item $X=\{b^k\mid k\in\mathbb{N}\}$ è $b$-automatico.
    \item $X=\{x\in\mathbb{N}\mid [x]_2 \text{ contiene un numero dispari di } 1\}$ è $2$-automatico.
    \item (facoltativo) $X=\{2^k\mid k\in\mathbb{N}\}$ non è $3$-automatico. \textcolor{red}{???}
\end{enumerate}
\end{ex}

\begin{sol}\
\begin{enumerate}
    \item Dobbiamo mostrare che $[\mathbb{N}]_b=\{[x]_b\mid x\in\mathbb{N}\}$ è DFA-decidibile.
    Basta controllare che la parola non inizi con ``zero'' e che, se lo fa, debba essere
    la parola ``zero''. Il DFA è:
    \[
    \begin{tikzpicture}[
        >=stealth', shorten >=1pt, auto,
        every state/.style={circle, draw, fill=white, minimum size=0.8cm, inner sep=1pt}
    ]
    % Stati
    \node[state]            (q0) at (0,0)    {$q_0$};
    \node[state, accepting] (q1) at (4,1.5)  {$q_1$};
    \node[state, accepting] (q2) at (4,-1.5) {$q_2$};
    \node[state]            (q3) at (8,1.5)  {$q_3$};

    % Freccia d'ingresso
    \draw[->] ([shift={(-0.7,0.7)}]q0.center) -- (q0);

    % Transizioni
    \draw[->] (q0) -- node {$0$} (q1);
    \draw[->] (q0) -- node[swap] {$i\in\{1,\dots,b-1\}$} (q2);
    \draw[->] (q1) -- node {$i\in\{0,\dots,b-1\}$} (q3);

    % Cappi
    \path[->] (q2) edge[loop below] node {$i\in\{0,\dots,b-1\}$} (q2);
    \path[->] (q3) edge[loop above] node {$i\in\{0,\dots,b-1\}$} (q3);
    \end{tikzpicture}
    \]

    \item Si ha $[X]_2=\{[x]_2\mid x=2^k,\ k\in\mathbb{N}\}=\{1\,0^k\mid k\in\mathbb{N}\}$. Quindi il DFA è:
    \[
    \begin{tikzpicture}[
        >=stealth', shorten >=1pt, auto,
        every state/.style={circle, draw, fill=white, minimum size=0.8cm, inner sep=1pt}
    ]
    % Stati
    \node[state]            (q0) at (0,0)    {$q_0$};
    \node[state, accepting] (q1) at (4,0)    {$q_1$};
    \node[state]            (q2) at (2,-2.5) {$q_2$};

    % Freccia d'ingresso
    \draw[->] ([shift={(-0.7,0.7)}]q0.center) -- (q0);

    % Transizioni
    \draw[->] (q0) -- node {$1$} (q1);
    \draw[->] (q0) -- node[swap] {$0$} (q2);
    \draw[->] (q1) -- node {$1$} (q2);

    % Cappi
    \path[->] (q1) edge[loop above] node {$0$} (q1);
    \path[->] (q2) edge[loop below] node {$0,1$} (q2);
    \end{tikzpicture}
    \]

    \item Si ha $[X]_b=\{[x]_b\mid x=b^k,\ k\in\mathbb{N}\}=\{1\,0^k\mid k\in\mathbb{N}\}$
    (il punto precedente è il caso $b=2$). Il DFA è:
    \[
    \begin{tikzpicture}[
        >=stealth', shorten >=1pt, auto,
        every state/.style={circle, draw, fill=white, minimum size=0.8cm, inner sep=1pt}
    ]
    % Stati
    \node[state]            (q0) at (0,0)    {$q_0$};
    \node[state, accepting] (q1) at (5,0)    {$q_1$};
    \node[state]            (q2) at (2.5,-2.5) {$q_2$};

    % Freccia d'ingresso
    \draw[->] ([shift={(-0.7,0.7)}]q0.center) -- (q0);

    % Transizioni
    \draw[->] (q0) -- node {$1$} (q1);
    \draw[->] (q0) -- node[swap] {$i\in\{0,2,\dots,b-1\}$} (q2);
    \draw[->] (q1) -- node {$i\in\{1,\dots,b-1\}$} (q2);

    % Cappi
    \path[->] (q1) edge[loop above] node {$0$} (q1);
    \path[->] (q2) edge[loop below] node {$i\in\{0,\dots,b-1\}$} (q2);
    \end{tikzpicture}
    \]

    \item Si ha $[X]_2=\{[x]_2\mid [x]_2 \text{ contiene un numero dispari di } 1\}$.
    Il DFA è (lo stato $q_z$ rifiuta le parole che iniziano con $0$):
    \[
    \begin{tikzpicture}[
        >=stealth', shorten >=1pt, auto,
        every state/.style={circle, draw, fill=white, minimum size=0.8cm, inner sep=1pt}
    ]
    % Stati
    \node[state]            (q0) at (0,0)    {$q_0$};
    \node[state, accepting] (qd) at (3.5,0)  {$q_d$};
    \node[state]            (qp) at (7,0)    {$q_p$};
    \node[state]            (qz) at (0,-2.5) {$q_z$};

    % Freccia d'ingresso
    \draw[->] ([shift={(-0.7,0.7)}]q0.center) -- (q0);

    % Transizioni
    \draw[->] (q0) -- node {$1$} (qd);
    \draw[->] (q0) -- node[swap] {$0$} (qz);
    \draw[->] (qd) to[bend left=25]  node {$1$} (qp);
    \draw[->] (qp) to[bend left=25]  node {$1$} (qd);

    % Cappi
    \path[->] (qd) edge[loop above] node {$0$} (qd);
    \path[->] (qp) edge[loop above] node {$0$} (qp);
    \path[->] (qz) edge[loop below] node {$0,1$} (qz);
    \end{tikzpicture}
    \]

    \item \textcolor{red}{???}
\end{enumerate}
\end{sol}

\newpage

\begin{ex}
Dopo aver ricordato la definizione della formula di esecuzione per un automa, dare la
matrice $M$ associata al seguente automa di alfabeto $A=\{a,b\}$ e calcolarne la serie
parziale $S$ troncata al terzo termine:
\[
S=M^0+M^1+M^2
\]
\[
\begin{tikzpicture}[
    >=stealth', shorten >=1pt, auto,
    every state/.style={circle, draw, fill=white, minimum size=0.8cm, inner sep=1pt}
]
% Stati
\node[state]            (q0) at (0,0)     {$q_0$};
\node[state]            (q3) at (3,1.5)   {$q_3$};
\node[state]            (q4) at (6,1.5)   {$q_4$};
\node[state, accepting] (q1) at (3,-1.5)  {$q_1$};
\node[state, accepting] (q2) at (6,-1.5)  {$q_2$};

% Freccia d'ingresso
\draw[->] ([shift={(-0.7,0.7)}]q0.center) -- (q0);

% Transizioni
\draw[->] (q0) -- node {$\varepsilon$} (q3);
\draw[->] (q0) -- node[swap] {$\varepsilon$} (q1);
\draw[->] (q1) -- node {$b$} (q3);
\draw[->] (q3) -- node {$a$} (q4);
\draw[->] (q4) -- node[pos=0.3] {$b$} (q1);
\draw[->] (q1) -- node[swap] {$a$} (q2);
\draw[->] (q2) -- node {$b$} (q4);

% Cappi
\path[->] (q4) edge[loop above] node {$b$} (q4);
\path[->] (q1) edge[loop below] node {$a$} (q1);
\end{tikzpicture}
\]
\end{ex}
\begin{sol}
Sia $\mathcal{A}$ un automa con $n$ stati, $M$ la sua matrice di adiacenza (a coefficienti in $\mathbb{N}\langle\langle A^*\rangle\rangle$, con $\bot$ per l'assenza di transizione, cioè la serie nulla), $I$ il vettore riga degli stati iniziali e $F$ il vettore colonna degli stati finali. La \emph{formula di esecuzione} è
\[
S_{\mathcal{A}}=I\cdot M^*\cdot F=I\Bigl(\sum_{k\ge 0}M^k\Bigr)F,
\]
dove il coefficiente della parola $w$ in $S_{\mathcal{A}}$ conta i cammini accettanti di etichetta $w$. Le $\varepsilon$-mosse si inseriscono in $M$ con il simbolo $\varepsilon$, interpretato come la serie formale della parola vuota.

\begin{enumerate}
\item Con gli stati ordinati $q_0,\dots,q_4$ la matrice associata è
\[
M=
\begin{array}{c|ccccc}
 & q_0 & q_1 & q_2 & q_3 & q_4\\
\hline
q_0 & \bot & \varepsilon & \bot & \varepsilon & \bot\\
q_1 & \bot & a & a & b & \bot\\
q_2 & \bot & \bot & \bot & \bot & b\\
q_3 & \bot & \bot & \bot & \bot & a\\
q_4 & \bot & b & \bot & \bot & b
\end{array}
\]
mentre
\[
I=\begin{pmatrix}\varepsilon&\bot&\bot&\bot&\bot\end{pmatrix},
\qquad
F=\begin{pmatrix}\bot&\varepsilon&\varepsilon&\bot&\bot\end{pmatrix}^{T}.
\]

\item Calcoliamo le potenze. $M^0$ è la matrice identità, con $\varepsilon$ sulla diagonale:
\[
M^0=
\begin{pmatrix}
\varepsilon & \bot & \bot & \bot & \bot\\
\bot & \varepsilon & \bot & \bot & \bot\\
\bot & \bot & \varepsilon & \bot & \bot\\
\bot & \bot & \bot & \varepsilon & \bot\\
\bot & \bot & \bot & \bot & \varepsilon
\end{pmatrix},
\]
\[
M^2=
\begin{pmatrix}
\bot & a & a & b & a\\
\bot & aa & aa & ab & ab+ba\\
\bot & bb & \bot & \bot & bb\\
\bot & ab & \bot & \bot & ab\\
\bot & ba+bb & ba & bb & bb
\end{pmatrix}.
\]

\item Sommando le tre matrici si ottiene
\[
S=M^0+M^1+M^2=
\begin{pmatrix}
\varepsilon & \varepsilon+a & a & \varepsilon+b & a\\
\bot & \varepsilon+a+aa & a+aa & b+ab & ab+ba\\
\bot & bb & \varepsilon & \bot & b+bb\\
\bot & ab & \bot & \varepsilon & a+ab\\
\bot & b+ba+bb & ba & bb & \varepsilon+b+bb
\end{pmatrix}.
\]
\end{enumerate}
\end{sol}

\newpage

\begin{ex}
Si consideri una macchina di Turing $\mu$ di alfabeto $A=\{0,1\}$ e sia
\[
X=\{x\in A^*\mid \text{la computazione di }\mu\text{ associata a }x\,x^R\text{ termina}\}
\]
dove $x^R$ è la parola $x$ invertita.

Mostrare che $X$ è un insieme semidecidibile definendo esplicitamente la macchina $\mu'$
che semidecide $X$ in funzione della macchina $\mu$.
\end{ex}

\begin{sol}
Fissata $\mu$, sia $Y=\{y\in A^*\mid \text{la computazione di }\mu\text{ associata a }y\text{ termina}\}$,
così che $X=\{x\in A^*\mid x\,x^R\in Y\}$. Sia $\alpha$ la macchina a un nastro che
semidecide $Y$ (ad esempio $\mu$ stessa).

Definiamo $\gamma$ come la macchina di Turing a $2$ nastri che, sul nastro $1$ legge $w$ e
scrive $w\,w^R$ sul nastro $2$. Le sue transizioni, con $b\in\{0,1\}$ e nella forma
$(\text{nastro }1,\text{nastro }2)/(\text{nastro }1,\text{nastro }2),(\text{mossa }1,\text{mossa }2)$, sono:
\[
\begin{tikzpicture}[
    scale=0.75, transform shape,
    >=stealth', shorten >=1pt, auto,
    every state/.style={circle, draw, fill=white, minimum size=0.8cm, inner sep=1pt}
]
% Stati
\node[state] (q0) at (0,0)   {$q_0$};
\node[state] (q1) at (7,0)   {$q_1$};
\node[state] (q2) at (14,0)  {$q_2$};
\node[state] (qa) at (21,0)  {$q_\alpha$};

% Freccia d'ingresso
\draw[->] ([shift={(-0.7,0.7)}]q0.center) -- (q0);

% Copia di w sul nastro 2
\path[->] (q0) edge[loop above, looseness=8] node {$(b,\square)/(b,b),(+1,+1)$} (q0);
\draw[->] (q0) -- node[swap] {$(\square,\square)/(\square,\square),(-1,0)$} (q1);

% Scrittura di w^R sul nastro 2
\path[->] (q1) edge[loop above, looseness=8] node {$(b,\square)/(b,b),(-1,+1)$} (q1);
\draw[->] (q1) -- node[swap] {$(\square,\square)/(\square,\square),(0,-1)$} (q2);

% Ritorno all'inizio del nastro 2
\path[->] (q2) edge[loop above, looseness=8] node {$(\square,b)/(\square,b),(0,-1)$} (q2);
\draw[->] (q2) -- node[swap] {$(\square,\square)/(\square,\square),(0,+1)$} (qa);
\end{tikzpicture}
\]
\begin{enumerate}
    \item In $q_0$ la macchina copia $w$ sul nastro $2$, spostando entrambe le testine a destra.
    \item In $q_1$ la testina del nastro $1$ torna indietro leggendo $w$ al contrario, mentre sul
    nastro $2$ si scrive $w^R$ di seguito a $w$.
    \item In $q_2$ la testina del nastro $2$ torna all'inizio di $w\,w^R$ e si passa a $q_\alpha$, lo
    stato iniziale di $\alpha$.
\end{enumerate}

La macchina $\mu'$ opererà prima come $\gamma$ e poi, dopo che $\gamma$ ha finito, opererà sul
nastro $2$ come $\alpha$, mentre il nastro $1$ è fermo (ogni transizione di $\alpha$ viene
estesa lasciando fissa la testina del nastro $1$). La macchina $\mu'$ si ferma solamente
quando si ferma $\alpha$.

Quindi $\mu'$ termina su $w$ se e solo se $\alpha$ termina su $w\,w^R$, cioè se e solo se
$w\,w^R\in Y$, ovvero $w\in X$. Dunque $\mu'$ semidecide $X$.
\end{sol}